\begin{figure}[tb]
\centering
\begin{tikzpicture}[x=1cm,y=1cm,>=Latex,font=\small]
 \fill[blue!3] (-3,-1.55) rectangle (3,1.55);
 \draw[very thick,blue!65!black] (-3,-1.55)--(-3,1.55);
 \draw[very thick,blue!65!black] (3,-1.55)--(3,1.55);
 \draw[very thick,orange!75!black] (-3,-1.55)--(3,-1.55);
 \draw[very thick,orange!75!black] (-3,1.55)--(3,1.55);
 \draw[thick,->] (-3,-.35) .. controls (-1.8,.5) and (1.8,-.5) .. (3,.35);
 \draw[thick,->,purple!75!black] (-.45,-1.55) .. controls (.8,-.5) and (-.8,.5) .. (.45,1.55);
 \fill (0,0) circle (2pt);
 \node[above right] at (.15,.1) {intersection};
 \node[below] at (-1.9,-.03) {$j(N)$};
 \node[left,purple!75!black] at (-.45,-1) {$A_+$};
 \node[left,blue!65!black] at (-3,0) {$E_R$};
 \node[right,blue!65!black] at (3,0) {$E_R$};
 \node[above,orange!75!black] at (1.9,1.55) {$E_B$};
 \node[below,orange!75!black] at (1.9,-1.55) {$E_B$};
 \node at (-2.1,1.05) {$P$};
\end{tikzpicture}
\caption{A schematic transverse section of the compact triad.
The relative cycle $j(N)$ has boundary on the radial side $E_R$ and avoids
$E_B$. Its dual representative $A_+$ has boundary on $E_B$ and is pushed
off $E_R$. Their nonzero intersection becomes a local degree after the
parameter is pulled back and doubled. The drawing records incidence and
boundary conditions, not the dimensions of the cycles.}
\label{fig:degree-triad}
\end{figure}
